\section{Weak linkage identification is a financial inverse problem}
\label{s61:linkage-sets}
The inability to distinguish two candidate histories is structurally related
to a weak instrument. In Equation~\eqref{s6:pair-mixing}, the eigenvalues of
the assignment matrix are $1$ and $r_s=2p_s-1$. The portfolio sum survives
even when the difference is unidentified. A Wilson interval containing
$1/2$ usefully stops the regular inverse, but that screening rule does not
construct an identification-robust confidence set. Nor does a large number
of production dates create audit information. The relevant analogy with
Section~\ref{r3:weak-iv} is to build uncertainty for observable quantities
before dividing by an uncertain identification coefficient
\citep[Section 2D, Theorem 3]{staiger1994}.

Use the fixed financial coefficient vector $k$ from
Equation~\eqref{s6:funding-account}, including the cost of paying the higher
first-month rate on induced balances. Within stratum $s$, define
\[
 m_{s1}=k'\eta_{s1},\quad m_{s2}=k'\eta_{s2},\quad
 S_s=\frac{m_{s1}+m_{s2}}2,\quad D_s=\frac{m_{s1}-m_{s2}}2.
\]
Let $v_{sA}=k'\theta_{sA}$ and $v_{sB}=k'\theta_{sB}$. Adding and
subtracting the two mixing equations gives
\begin{equation}
 S_s=\frac{v_{sA}+v_{sB}}2,\qquad
 D_s=r_s\frac{v_{sA}-v_{sB}}2=r_s(v_{sA}-S_s).
 \label{s61:diagonal-map}
\end{equation}
Thus $v_{sA}=S_s+D_s/r_s$ for $r_s\ne0$. At $r_s=0$, the observable
restriction in these mean equations is $D_s=0$ and $v_{sA}$ is unrestricted
by them because $v_{sB}$ can compensate. The final funding value is
$V=\sum_s w_s v_{sA}-C$, with fixed $w_s=1/2$ and
$C=Nc r B_0$. These are monetary contrasts before the common fixed price
cost is subtracted; no balance stock is reinterpreted as a flow.

This is an inverse problem for response moments. It is not a claim that
every value is compatible with every feature of the full observed mixture
distribution. At $p=1/2$, conditional on an action/history, that distribution
has the symmetric form
$\{P(Y_A,Y_B)+P(Y_B,Y_A)\}/2$. Exchanging the two named portfolio laws
preserves it, so focal identity can remain ambiguous. Additional distributional
restrictions may nevertheless identify unordered components or their
separation and produce a smaller set than the mean equations permit.
The procedure below constructs a conservative outer confidence set without
using those restrictions; it does not compute the sharp identified set of
a fully specified mixture. Its independent audit supplies correctness counts
only. Audited action/outcome histories could directly inform portfolio
responses and would define a richer experiment.

\subsection{A confidence region for regular observable information}
The candidate-response vector is estimated before inverse linkage. Its
calendar influence and HAC covariance are those in
Section~\ref{s6:linked-histories}; every financial reduction uses the entire
within-horizon covariance. Apply the fixed linear transformation above to
obtain $\widehat z=(\widehat S_0,\widehat D_0,
\widehat S_1,\widehat D_1)$ and its covariance $\widehat V_z$.
For total error $\alpha$, allocate $\alpha/2$ to the four observable
coordinates. Marginal bands are
\begin{equation}
 I_j=\left[\widehat z_j-z_{1-\alpha/16}\sqrt{\widehat V_{z,jj}},
                 \widehat z_j+z_{1-\alpha/16}\sqrt{\widehat V_{z,jj}}\right].
 \label{s61:primitive-bands}
\end{equation}
Each uses error $\alpha/8$, giving joint coverage at least
$1-\alpha/2+o(1)$ by Bonferroni when its studentized approximation is valid.
There is no independence assumption between these four estimates.

The approximation concerns candidate means, not a near-singular causal
inverse. It retains the independent pilot, finite correct basis, fixed
horizon, nonvanishing stratum/innovation masses, mixing/moment conditions
and consistent lag covariance established for the funding law. The
bandwidth conditions of \citet[Theorem 2]{newey1986} remain relevant after a
linear transform. Near-unit-root dynamics, unqualified nuisance learning or
a missing evaluation stratum are not repaired by projecting a confidence
region. Positive limiting variances are required for the stated normal
bands; a deterministically known coordinate instead needs an exact
singleton region justified by its model.

For each independent representative audit let
$K_s\sim\operatorname{Binomial}(M_s,p_s)$. Allocate error
$\epsilon=\alpha/4$ to its two-sided probability interval $P_s=[L_s,U_s]$.
For $0<K_s<M_s$, obtain the endpoints from
\begin{equation}
 \begin{split}
 \sum_{j=K_s}^{M_s}\binom{M_s}{j}L_s^j(1-L_s)^{M_s-j}&=\epsilon/2,\\
 \sum_{j=0}^{K_s}\binom{M_s}{j}U_s^j(1-U_s)^{M_s-j}&=\epsilon/2.
 \end{split}
 \label{s61:binomial-inversion}
\end{equation}
Set $L_s=0$ at zero successes, $U_s=1$ at all successes, and
$P_s=[0,1]$ if $M_s=0$. The binomial tail probabilities are monotone
in $p$. For either endpoint the rejecting counts form one tail whose
total probability is bounded by $\epsilon/2$. Taking the union of the
two rejection events proves
$P_p\{p\notin P_s\}\le\epsilon$ for every $p\in[0,1]$.
Discreteness can make coverage conservative. This proof does not require
$M_s$ to grow with production dates; a fixed audit simply leaves a region
that need not shrink. These claims require accurate, noninformatively
sampled audit labels. The audit assumptions in
\citet[Section 2.1, Audit Sample]{slawski2025} make that distinction explicit;
they are not consequences of the interval formula.

\subsection{Projection, coverage and the shape of the result}
Let $\mathcal R$ be the Cartesian product of the four observable bands and
the two audit intervals. For each stratum construct
\begin{align}
 \mathcal V_s&=\{v:\exists S\in I_{S_s},D\in I_{D_s},p\in P_s,
                         \ D=(2p-1)(v-S)\},\nonumber\\
 \mathcal V&=\left(\sum_s w_s\mathcal V_s\right)-C.
 \label{s61:projection}
\end{align}
The sum denotes all sums of one element from each set, not the sum of
arbitrarily selected endpoints. On the event that all primitive quantities
lie in $\mathcal R$, the true $v_{sA}$ satisfy the displayed equations,
including at $p_s=1/2$. Hence the true $V$ lies in $\mathcal V$ on that
same event. The union bound gives
\[
 P\{V\in\mathcal V\}\ \ge\ 1-\alpha+o(1).
\]
This is a project coverage argument, not a new theorem attributed to the
linkage paper or a direct application of the Gaussian Anderson--Rubin test.
The audit region is finite-sample valid; the response region remains
asymptotic. Projection itself needs no lower bound on $|2p_s-1|$.
A uniform claim over a sequence of linkage probabilities additionally
requires uniform validity of the primitive response bands along that
sequence. Pointwise CLTs cannot establish that claim by renaming the result.
The distinction between a finite normal experiment and an asymptotic
approximation is also explicit in
\citet[Section 2.5, equations (9)--(11)]{rambachan2023}.

There is a complete scalar algorithm without a bounded effect grid.
Write $I_S=[a,b]$, $I_D=[c,d]$ and
$R=[2L-1,2U-1]$. On each positive or negative part of $R$,
the extrema of $D/r$ are its four rectangle corners, with signed infinite
limits at zero. Add $[a,b]$, retain both sign components, and merge only
overlapping closed intervals. If $0\in R$ and $0\in I_D$, the rank-one
equation permits every real $v$, so return $\mathbb R$. If $R=\{0\}$
and $0\notin I_D$, return the empty set: that primitive region is
incompatible with the model. Weighted interval sums then produce
Equation~\eqref{s61:projection}. This finite interval arithmetic preserves
unbounded rays and disconnected components.

For example, $S=1.5$, $D=.5$ and $p\in[.4,.6]$ give
\[
 \mathcal V_s=(-\infty,-1]\ \cup\ [4,\infty).
\]
Neither a finite interval joining plotted points nor a selected positive
branch describes this information. As an independent check on the algorithm,
fix a candidate $v$ and ask whether the rectangle image
$\{r(v-S):r\in R,S\in[a,b]\}$ intersects $[c,d]$.
Its endpoints are four products, so this test uses no ratio. The executed
16,800 comparisons agree, including exactly singular regions, and a finite
enumeration over all binomial counts for sizes 1--25 and 201 probability
values gives minimum audit coverage .95096 for nominal .95. These checks
validate finite identities; they are not a proof of the calendar CLT.

\subsection{Precision, alternatives and the bank's next decision}
A second valid projection uses separate bands for $(m_{s1},m_{s2})$ and
the original inverse map. It ignores much of their correlation when forming
a rectangle. The sum/difference construction uses that covariance before
forming bands: a common monetary shock cancels from the difference, while
uncertainty in the sum need not be amplified by $1/|2p-1|$.
Neither rectangular construction is uniformly shorter for every covariance
law. The registered comparison therefore evaluates both on the same draws,
alongside the regular delta interval and its availability rule.

A joint ellipsoidal response region could exploit further covariance but
requires reliable constrained projection, including unbounded directions.
Alternatively, for proposed $(v,p)$ one can test
$\widehat D-(2p-1)(v-\widehat S)$ using its candidate-specific covariance
and combine the audit evidence before eliminating $p$. Its variance includes
$V_{DD}+(2p-1)^2V_{SS}+2(2p-1)V_{SD}$; using a fixed variance at a fitted
inverse would recreate the weak-normalization error. A justified joint
test could improve power, but its critical value, nuisance elimination and
degenerate cases require a separate derivation. The selected rectangle
projection buys transparent coverage and exact topology at the cost of
conservatism. A joint Bayesian linkage/outcome model
\citep{guha2022} instead uses a different likelihood, feedback and prior
restriction; its posterior is not this confidence set.

The five-regime funding comparison uses 720 evaluation dates, an independent
2,880-date nuisance pilot and 200 independent audits per stratum, with 400
repetitions in each regime. Its true financial value is about 375,149 currency
units. The response law, target and error allocation remain fixed while the
mapping probabilities change.

\begin{table}[htbp]\centering\small
\caption{Weak linkage and financial inference. Projection width is the median
conditional on a bounded set, in millions of currency units. Regular coverage
counts unavailable intervals as no coverage.}
\label{s61:projection-results}
\begin{tabular}{lrrrrr}\toprule
$(p_0,p_1)$ & Regular & Regular & Projection & Bounded & Projection\\
 & available & coverage & coverage & fraction & width\\\midrule
(0.95, 0.80) & 400/400 & 0.9250 & 1.000 & 1.0000 & 1.498 \\
(0.70, 0.65) & 397/400 & 0.9150 & 1.000 & 0.9475 & 2.947 \\
(0.55, 0.52) & 15/400 & 0.0350 & 1.000 & 0.0000 & -- \\
(0.50, 0.50) & 2/400 & 0.0025 & 1.000 & 0.0000 & -- \\
(0.20, 0.35) & 395/400 & 0.9225 & 1.000 & 0.9550 & 2.542 \\
\bottomrule\end{tabular}
\end{table}

Projection covers the retained truth in all 400 draws in every regime; a
95\% binomial Monte Carlo interval for each coverage rate is $[.9908,1]$.
That result is conservative, not evidence of exact 100\% coverage. The weak
and singular cases return the whole real line in every draw. The moderate
case contains one disconnected result and twenty whole-line results; the
reversed case has eighteen whole-line results. Conditional median widths
for the candidate-mean rectangle are 1.910, 6.644 and 5.197 million in the
strong, moderate and reversed regimes, versus 1.498, 2.947 and 2.542 million
for sum/difference projection. Their conditional comparison excludes
unbounded sets; it cannot be interpreted as an unconditional width guarantee.

The regular procedure is substantially narrower where available, but it
covers .925 of the strong-regime draws, .915 of all moderate requests and
.9225 of all reversed requests. In the last two regimes its coverage
conditional on availability is .9219 and .9342. Even the strong-regime
finding is an adverse finite-sample result alongside the earlier
100/400-audit study; a new independent seed set does not have to reproduce
its favorable coverage. In the weak case the regular rule produces only
15 intervals, 14 covering truth; in the singular case it produces two,
only one covering truth. Those conditional denominators cannot justify
promoting the screened interval to inference robust to weak linkage.


An unbounded set can be the correct answer to the supplied information.
If an operationally justified lower bound on mapping accuracy excludes
$1/2$, it changes the admissible model and may recover a bounded result.
More representative correctness-count audits can resolve a genuinely nonzero
$r_s$, but cannot invert these mean equations when $p_s$ is exactly $1/2$.
In that case stronger identifiers, an additional measurement, or a changed
economic target such as the portfolio sum is required. Mislabelled audits,
selected audits and a persistent crosswalk outside the declared per-origin
map law require different observation models. High coverage obtained by
returning $\mathbb R$ is honest lack of information, not a useful investment
case or evidence of universal causal inference.
